\section{Independent scope and reproducibility audit}
\label{sec:p19-audit}
The research programme combines a stipulated inverse-time relay progress model, an exact independent-rectangle uncertainty criterion under fixed scalar-scaled current templates, a correlated-derating interior counterexample, synthetic comparative screening and model-level IEEE 123-bus feeder experiments. These components have different evidentiary status and must not be conflated. In particular, the P13 rectangular result is an if-and-only-if statement only for Cartesian current-scale and breaker-delay uncertainty with fixed traces, fixed reset dynamics and finite required assertions. Backup clearing is evaluated under a sustained-fault counterfactual; the backup assertion time, rather than its breaker clearing time, enters the stipulated primary-to-backup coordination margin.

An independent deterministic replay of the \(90\) seeded synthetic scenarios recovered \(22\) complete-model approvals. Against that model reference, nominal, common-endpoint, independent-rectangle and finite time-domain-grid screens produced \(19,3,0,0\) false approvals. The seven P18 single-component ablations produced \(14,1,6,22,0,3,0\) false approvals in the order frequency, primary deadline, backup deadline, coordination margin, timer-memory approximation, independent derating and breaker-delay omission. The timer-memory approximation caused two conservative rejections. These are model-relative confusion counts, not independently observed field protection outcomes. The finite time-domain grid shares the stipulated relay integrator with the reference and therefore is not an external numerical oracle.

The P13 analytic common-derating counterexample has a coordination margin above \(0.4\) s at the two uncertainty endpoints but below \(0.4\) s at an interior value. This demonstrates failure of common-diagonal endpoint sufficiency without contradicting the exact independent-rectangle criterion. The P14--P16 feeder experiments extract source-pinned model information and show that disabling the documented \texttt{Line.sw1} element yields zero reported voltage at the monitored buses, even when the solver converges. These observations do not establish a physically feasible restoration sequence, independently verified transformer/regulator phase transfer, energized conductor continuity, installed relay zones, CT orientation, or field thermal adequacy.

\paragraph{Submission readiness.}
At the time of this audit, the repository's \texttt{shes/main.tex} remained a placeholder and the full bibliography and four manuscript figure PDFs were absent. Thus a complete manuscript build from a clean GitHub checkout, anonymous PDF inspection, publisher-level bibliography verification, and figure regeneration could not be certified. The separate P13--P19 sections are research artefacts pending incorporation into the complete manuscript. No physical restoration authorization is asserted.